\begin{pastexam}[label={px:ncku-107-ee-ds-02}]{成大電機所 107 學年度 複選題第 3 題}
（複選；每題 5 分，全對才給分）

Which of the following statement(s) is (are) true?
\begin{enumerate}[label=(\Alph*)]
\item Quick sort and merge sort both apply the divide-and-conquer approach to solve the sorting
problem.
\item Quick sort spends large effort in dividing a sequence into two parts while merge sort spends
large effort in combining two unsorted sequence into one sorted sequence.
\item Both of the two algorithms are stable.
\item The expected time complexity of merge sort is O(nlogn) which is faster than quick sort in the
worse case.
\end{enumerate}
\end{pastexam}

\begin{pastsoln}
答案：(A)(B)(D)。

(B) 的「two unsorted sequence」用字不精確：合併排序在合併那一步拿到的兩段，已經由遞迴排好了。 這裡
把它讀成「把原本未排序的資料分成的兩段，合併成一個排好的序列」，(B) 要表達的是「快速排序的功夫花在
分、合併排序的功夫花在合」，這是對的。 若硬要照字面說「合併兩個未排序的序列」，(B) 就是錯的，答案
變成 (A)(D)；下面照出題的意思作答。

\textbf{(A) 對}。 兩者都是分治法（資結書第~14.4、14.5~節）：把問題分成子問題、遞迴解、再組合。 快速排序
以基準分割成「小於等於基準」和「大於基準」兩段，分別遞迴；合併排序從中間切成兩半，分別遞迴後合併。

\textbf{(B) 對（照上面的讀法）}。 快速排序的工作都在分割：分割完兩段各自排好，整個陣列就排好了，組合
不必做任何事。 合併排序剛好相反：切一半是 $O(1)$ 的，工作都在合併兩個排好的半段，每層合併共 $\Theta(n)$。

\textbf{(C) 錯}。 合併排序是穩定的（合併時兩邊鍵值相等，先取左半段的）；快速排序不穩定（資結書定義~14.1）。
反例：用 Lomuto 分割（以最後一個為基準）排 $\langle 2_a, 2_b, 1 \rangle$，基準是 $1$，$2_a$、$2_b$ 都比它大，
最後基準和 $A[1] = 2_a$ 交換，得到 $\langle 1, 2_b, 2_a \rangle$，兩個 $2$ 的相對順序反了。

\textbf{(D) 對}。 合併排序不論輸入如何都是 $\Theta(n \lg n)$，期望時間當然也是 $O(n \lg n)$；快速排序的最壞
情況（例如已排好的輸入、每次以最後一個為基準）分割成 $n - 1$ 和 $0$，是 $\Theta(n^2)$。 所以就最壞情況而言，
合併排序比快速排序快。 （兩者的期望時間都是 $\Theta(n \lg n)$，快速排序常數小，實務上通常比較快。）
\end{pastsoln}
